% Last Updated: March 28, 2023
	\documentclass[a4paper, twoside]{article}
	% Margins
	\usepackage{geometry}
		\geometry{
			% showframe,
			left   = 20mm,
			right  = 20mm,
			top    = 20mm,
			bottom = 20mm
		}

	% Headers and footers
	\usepackage{fancyhdr}
		\pagestyle{fancy}
		\fancyhf{}
		\fancyhead[LE, RO]{\thepage}
		\fancyhead[CE]{\textsc{Cong Wen}}
		\fancyhead[CO]{\textsc\rightmark}

		\setlength{\headheight}{15pt}
		\renewcommand\headrulewidth{0pt}








\input{/Users/wenc/Documents/Math/universal-pre}

\title{\tbf{Faltings height of abelian varieties}}
\author{Cong Wen, 05/2023}
\date{}
% Last Updated: March 28, 2023



\begin{document}
\maketitle
\thispagestyle{empty}
\begin{abstract}
	This is my self-use notes\footnote{Last updated on October 19, 2023} to understand arithmetic topics around \tbf{Faltings height of abelian varieties}.

	Please reach out if you find anything incorrect.

\end{abstract}
\tableofcontents
% ----------------------------------------------------------------------------------------------------

% \sct*{Principles}
% \begin{itm}
% 	\item Collect fundamentals and related properties among literatures.
% 	\item Learn those known methods tried by others among literatures.
% 	\item Identify the essential difficulties, which should be aligned into a tree.
% 	\item Try to overcome one or few leaf nodes of the essential-difficulty tree.
% 	\item Avoid generalizing to more technical terms when doing so only blurs things.
% \end{itm}


\sct{Outline}

\prg{Ultimate goal} Develope whatever is necessary to effectively compute the Faltings height of an abelian variety $A$ over any number field $K$ or find their relationship with other analytic objects, study the effect on Hodge bundles along $f$, study the change in Faltings height along any morphism $f:A\ar B$, including non isogenies.


\prg{Some heuristics} For other heuristics and relations, useful or not.
\begin{itm}
	\item A series of papers of Ullmo, Szpiro-Ullmo, Abbes-Ullmo in 2000s proved many results in the $J_{0}(N)$ case, as well as quotient abelian varieties of them which are not isogenies! However all those assume $N$ is square-free, and the constants involved is not effective.
	\item Hyperbolicity of moduli spaces is a general phenomenon, the theory of Hodge modules developed by Saito, Popa-Schnell et al. plays an important role. One key along this line also features the study of the singular metrics on the Hodge bundle.
	\item Odaka has some heuristics relating Faltings height to Futaki invariants.
	\item It is important to establish the modularity of the generating series of special cycles in many cases, but can one find a converse theorem analogous to the automorphic-to-Galois direction, i.e., for any modular forms with some conditions, there must exist some special cycles whose generating series match with it?
	\item The Hodge bundle give rise to the Faltings height, which is a logarithmic derivative in CM case. The theta divisor give rise to the N{\'e}ron-Tate height, which is a derivative for quadratic points. Why these two? What might happen if I study other distinguished divisors on some moduli space?
\end{itm}


\prg{Formulae and bounds}
\begin{thm}[Bost1996]
	There exists an effective constant $C>0$ such that $h_{\text{Fal}}^{\text{st}}(A)\gqs-gC$.
\end{thm}


\begin{thm}[{\cite[Theorem~1.3]{Wei2023}}]
	Let $A$ be a CM abelian varietie, assume the Colmez conjecture holds for $A$, and no $\frac{1}{O(1)}$-Siegel zero conjecture of $L(s,\chi_{d})$ holds. Then there exists effective constant $c_{1}>0,c_{2}$ depending only on $g$ such that
	\[
		h_{\text{Fal}}^{\text{st}}(A)\lqs c_{1}\frac{\log\sP{\Dta_{K}}}{[K:\bQ]}+c_{2}.
	\]
\end{thm}

Of course, for CM elliptic curves, Faltings height should already have subpolynomial bounds in discriminants (Pila--Shankar--Tsimerman).

\sct{Fundamentals}

\begin{dfn}
	Let $K$ be a number field, $A\ar\Spec K$ be an abelian variety of dimension $g$, $\pi:\cA\ar\Spec\cO_{K}$ its N{\'e}ron model and $e:\Spec\cO_{K}\ar\cA$ the zero section. Take $\oga_{\cA}=e\xS \Oga_{\cA/\Spec\cO_{K}}^{g}$ and $\afa\in \mathrm{H}_{}^{0}(\Spec\cO_{K},\oga_{\cA})$, then the unstable Faltings height of $A$ is given by
	\[
		h_{\mathrm{Fal}}^{\mathrm{un}}(A/K)=\frac{1}{[K:\bQ]}\sR{\log\#(\oga_{\cA}/\afa\cO_{K})-\frac{1}{2}\sum_{\sgm:K\ar\bC}\log\sP{\int_{A_{\sgm}(\bC)}\afa\wge\oL{\afa}}}+C,
	\]
	and $h_{\mathrm{Fal}}^{\mathrm{st}}(A)=h_{\mathrm{Fal}}^{\mathrm{un}}(A_{K'}/K')$ for any finite extension $K'/K$ such that $A_{K'}$ is semistable.
\end{dfn}
\begin{rmk}
	For the normalization factor $C$,
	\begin{itm}
		\item \cite{Faltings1983,Silverman1986,SU1999,Ullmo2000} used $C=(g/2)\log{2}$.

		This makes height of elliptic curves neat.
		\item \cite{Yang2010Colmez,Yang2010Hilbert,YZ2018,YZ2021,Yuan2021Bigness} used $C=(g/2)\log{2\pi}$.

		This makes Colmez conjecture neat.
		\item \cite{Colmez1993,Colmez1998Hauteur,AGHMP2018} used $C=0$.

		This makes the original definition neat.
	\end{itm}
\end{rmk}

\begin{rmk}
	One can keep only the infinite part if the N{\'e}ron differential $\afa$ exists, in particular this is true after passing to a sufficiently large extension of $K$.
\end{rmk}

\begin{qst}
	Is there sufficient and necessary arithmetic geometry conditions that ensure two isogenous abelian varieties are isomorphic?

	If one specializes to CM elliptic curves, this more or less amounts to distinguishing different ideal classes. Thus a cheap criterion may not exist?
\end{qst}


\begin{thm}[{\cite[Proposition~1.1]{Silverman1986}}]
	Let $E$ be an elliptic curve over $\bQ$ with minimal discriminant $\Dta_{E}$ and let $E_{\bC}\cong\bC/(\bZ+\bZ\tau)$ with $\Img\tau>0$, then
	\[
		h_{\mathrm{Fal}}^{\mathrm{un}}(E/\bQ)=\frac{1}{12}\log\sP{\Dta_{E}}-\frac{1}{12}\log\sP{\Dta(\tau)}(\Img\tau)^{6}.
	\]
\end{thm}

\begin{thm}[{\cite[Proposition~3.1]{Silverman1986}}]
	Let $E$ be an optimal elliptic curve over $\bQ$, then only archimedean part contributes to the height, one has
	\[
		h_{\mathrm{Fal}}^{\mathrm{st}}(E/\bQ)=\frac{1}{2}\log\Deg\phi_{E}-\log\dP{f_{E}}-\log\sP{c_{E}}.
	\]
\end{thm}



\red{- - - under construction - - -}

\begin{thm}
	Let $A$ be an abelian variety over a number field $K$, then
	\begin{enr}[label = (\arabic*)]
		\item for any $\sgm\in \Gal(\oL{\bQ}/\bQ)$, one has
		\[
			h_{\mathrm{Fal}}^{\mathrm{un}}(A/K)=h_{\mathrm{Fal}}^{\mathrm{un}}(\sgm(A)/\sgm(K)),\quad h_{\mathrm{Fal}}^{\mathrm{st}}(A)=h_{\mathrm{Fal}}^{\mathrm{st}}(\sgm(A)).
		\]
		\item if there is an isomorphism $A\cong A_{1}\tms_{K} A_{2}$ over $K$, one has
		\[
			h_{\mathrm{Fal}}^{\mathrm{un}}(A/K)=h_{\mathrm{Fal}}^{\mathrm{un}}(A_{1}/K)+h_{\mathrm{Fal}}^{\mathrm{un}}(A_{2}/K),\quad h_{\mathrm{Fal}}^{\mathrm{st}}(A)=h_{\mathrm{Fal}}^{\mathrm{st}}(A_{1})+h_{\mathrm{Fal}}^{\mathrm{st}}(A_{2}).
		\]
		\item for any subfield $L\sbs K$ such that $K/L$ is Galois, let $B=\Res_{K/L}(A)$, one has
		\[
			h_{\mathrm{Fal}}^{\mathrm{un}}(B/L)=[K:L]h_{\mathrm{Fal}}^{\mathrm{un}}(A/K),\quad h_{\mathrm{Fal}}^{\mathrm{st}}(B)=[K:L]h_{\mathrm{Fal}}^{\mathrm{st}}(A).
		\]
	\end{enr}
\end{thm}

\begin{prf}
	The first claim is obvious because the finite part is already Galois invariant and the infinite part is a sum over all conjugates of $K$. For the second claim, extend the isomorphism to N{\'e}ron models $\cA\cong\cA_{1}\tms_{\cO_{K}}\cA_{2}$ over $\cO_{K}$, then $\Oga_{\cA/\cO_{K}}^{1}\cong\pr_{1}\xS \Oga_{\cA_{1}/\cO_{K}}^{1}\op\pr_{2}\xS \Oga_{\cA_{2}/\cO_{K}}^{1}$ and $\oga_{\cA}\cong\oga_{\cA_{1}}\ot_{\cO_{K}}\oga_{\cA_{2}}$ as $\cO_{K}$-modules of rank $1$, the additivity follows. The third claim directly follows from the first two and that $B\cong\prod_{\sgm\in\Gal(K/L)}\sgm(A)$ over $K$.
\end{prf}


\begin{qst}
	Is there a sufficient and necessary condition for the vanishing of the finite part of $h_{\mathrm{Fal}}^{\mathrm{st}}(A)$?
\end{qst}



\sct{Colmez product formula}

\ssc{Chowla--Selberg}

\begin{thm}[Chowla--Selberg]
	Let $K$ be an imaginary quadratic field of discriminant $-D$, class number $h$, unit group of order $w$, and Dirichlet character $\chi$. Let $E$ be an elliptic curve over $\bQ$ with complex multiplication by $\cO_{K}$. Then
	\begin{align*}
	  \frac{1}{12h}\sum_{\ka\in\Cl(\cO_{K})}\log(\Dta(\ka)\Dta(\ka\xI))&=\log(2\pi/D)+\frac{w}{2h}\sum_{k=0}^{D}\chi(k)\log\Gma(k/D), \\
	  h_{\mathrm{Fal}}^{\mathrm{st}}(E/\bQ)&=-\frac{1}{2}\frac{L'(0,\chi)}{L(0,\chi)}-\frac{1}{4}\log{D}-\frac{1}{2}\log{\pi}.
	\end{align*}

\end{thm}

\begin{prf}[Sketch of the proof]
	For the second identity, use the first identity and that for any $\ka\in\Cl(\cO_{K})$,
	\[
		\log{\sP{\Dta(\tau)}(\Img\tau)^{6}}=\frac{1}{2}\log{\Dta(\ka)\Dta(\ka\xI)}-3\log{(4/D)}.
	\]
	For the first identity, one can compare the analytic expansion of Kronecker and Lerch.

	\prg{Step 1} By Kronecker's limit formula, one has for each $\ka\in\Cl(\cO_{K})$,
	\[
	  \zta_{K}(s,\ka)=\sum_{[I]\in\ka}(\Nrm_{K/\bQ}(I))^{-s}=-\frac{1}{w}-\frac{1}{12w}\log(\Dta(\ka)\Dta(\ka\xI))s+O(s^{2}),
	\]
	therefore from $\zta_{K}(s)=\sum_{\ka\in\Cl(\cO_{K})}\zta_{K}(s,\ka)$, one has
	\[
		\zta_{K}(0)=\frac{h}{w},\quad \frac{\zta_{K}'(0)}{\zta_{K}(0)}=\frac{1}{12h}\sum_{\ka\in\Cl(\cO_{K})}\log{\Dta(\ka)\Dta(\ka\xI)}.
	\]

	\prg{Step 2} By Lerch's expansion for Hurwitz zeta function $H(s,x), 0<x\lqs1$, one has
	\[
		H(s,x)=\sum_{n=0}^{\ift}\frac{1}{(n+x)^{s}}=(1/2-x)+\log{\frac{\Gma(x)}{\sqrt{2\pi}}}s+O(s^{2}),
	\]
	therefore from $\zta_{\bQ}(s)=H(s,1), L(s,\chi)=D^{-s}\sum_{k=0}^{D}\chi(k)H(s,k/D)$, one has
	\[
		\zta_{\bQ}(0)L(0,\chi)=-\frac{1}{2}\sum_{k=0}^{D}\chi(k)(k/D),\quad \frac{\zta_{\bQ}'(0)}{\zta_{\bQ}(0)}+\frac{L'(0,\chi)}{L(0,\chi)}=\log(2\pi/D)-\frac{\sum_{k=0}^{D}\chi(k)\log\Gma(k/D)}{\sum_{k=0}^{D}\chi(k)(k/D)}.
  	\]

	\prg{Step 3} Now use $\zta_{K}(s)=\zta_{\bQ}(s)L(s,\chi)$ and compare the expansion of Kronecker and Lerch.
	\begin{itm}
		\item Evaluating at $s=0$ gives class number formula $\sum_{k=0}^{D}\chi(k)(k/D)=-2h/w$,
		\item Taking logarithmic derivatives at $s=0$ gives Chowla--Selberg formula.
	\end{itm}
\end{prf}
\begin{qst}
	One can see the key is the different treatment of the analytic objects in $\zta_{K}(s)=\zta_{\bQ}(s)L(s,\chi)$, what if one can extract higher coefficients and compare them?

\end{qst}

\ssc{Yuan--Zhang}

Use \cite{YZ2018,YZ2018Erratum}.

\begin{qst}
	Although it is already pointed out in \cite{YZZ2013} the impossibility of establishing the strongest $Z=\Phi$ equality for arbitrary Schwarz function, but what's the implications if one assumes it?
\end{qst}

\sct{Factorization estimate}

This is one possible technique to distinguish different isomorphism classes in an isogeny class. In particular, factorization estimate of Masser--W{\"u}stholz says one can use height to bound the isogeny degree.

\begin{thm}[{\cite[Theorem~I]{MW1995Factorization}}]
	For any $n,d\in\bZ_{>0}$, there exists constants $\kpa(n),C(n,d)$ satisfying that:

	If $A$ is an abelian variety of dimension $n$ over a number field $k$ of degree $d$, for any field extension $K$ of $k$, $A_{K}$ is isogenous to $A_{1}^{e_{1}}\tms\cdots\tms A_{t}^{e_{t}}$, where $\{A_{i}\}$ are mutually non-isogenous $K$-simple abelian subvarieties of $A_{K}$, with isogeny degree at most $C\max\{1,h(A)\}^{\kpa}$.
\end{thm}


\sct{Northcott property}

\begin{thm}
	Let $A$ be an abelian variety over a number field $K$, then
	\begin{enr}
		\item (Faltings) the $K$-isomorphism classes in the $K$-isogeny class of $A$ with bounded Faltings height is finite, moreover, they do have bounded Faltings height.
		\item (Kisin--Mocz) the $\oL{K}$-isomorphism classes in the $\oL{K}$-isogeny class of $A$ with bounded Faltings height is finite if the Mumford--Tate conjecture holds for $A$.
	\end{enr}
\end{thm}


\sct{Literature extraction}

\ssc{Faltings1983}


This part of the note is based on \cite{Faltings1983,Faltings1983Erratum,Paris1985}.

\begin{ups}
	According to Faltings' letter\cite{Paris1985} to Szpiro in which he announced the proof, one can see that the height theory of abelian varieties/Tate conjectures/Kodaira-Parshin construction/Northcott properties and many other keys are known inside out to him. The key moment is when Faltings found the argument to the Shafarevich conjecture, i.e., the isomorphism classes are finite within each isogeny class (he already knew that the isogeny classes are finite):
	\begin{itm}
		\item Fix some $A$ and consider those $B$ isogenous to $A$, by Northcott properties it suffices to show those $h_{\text{Fal}}(B)$ form a bounded set. First, any prime $\ell$, the $G_{K}$-lattice $T_{\ell}(B)\sbs V_{\ell}(B)\cong V_{\ell}(A)$ falls in finitely many isomorphism classes, so the $\ell$-adic valuation of $\exp(2[K:\bQ](h_{\text{Fal}}(B)-h_{\text{Fal}}(A)))$ is bounded. So it suffices to show for $\ell\gg0$, an isogeny $\phi:A\ar B$ of degree $\ell^{h}$ does not change the Faltings height. Assume $\ell\cdot\Ker\phi=0$, write $d=\#\Oga_{\Ker\phi}^{1}$ and $m=[K:\bQ]$, then $h_{\text{Fal}}(B)-h_{\text{Fal}}(A)=\log(\ell)(h/2-d/m)$, it suffices to prove $d=mh/2$ for $\ell\gg0$.
		\item Now consider the $G_{\bQ}$-modules,
		\[
			\Lda^{mh}(\Ind_{G_{K}}^{G_{\bQ}}(\Ker\phi(\oL{K})))\sbs\Lda^{mh}(\Ind_{G_{K}}^{G_{\bQ}}(A[\ell](\oL{K}))),
		\]
		the former $G_{\bQ}$-module is of dimension $1$, therefore $G_{\bQ}$ must act on the former as characters, and up to a power of sign $\eps:G_{\bQ}\ar\{\pm1\}$ it must act as a power of cyclotomic character $\chi:G_{\bQ}\ar(\bZ/\ell\bZ)\xT$ by class field theory. A theorem of Raynaud\cite[4.1.1]{Raynaud1974} says it is exactly $\chi^{d}$, therefore $\Frob_{p}$ acts as $\pm p^{d}$.
		\item Note that $P_{mh}(T)=\Det(T-\Frob_{p}|\Lda^{mh}(\Ind_{G_{K}}^{G_{\bQ}}(T_{\ell}(A))))$ is independent of $\ell$, and its root must have absolute value $p^{mh/2}$.
	Also note that $A[\ell]=T_{\ell}(A)/\ell T_{\ell}(A)$, therefore $\pm p^{d}$ must be a root of $P_{mh}(T)\bmod \ell$. One can take $\ell$ large enough, so $d$ must be $mh/2$. Therefore, the Shafarevich conjecture, as well as the Mordell conjecture, is proved.
	\end{itm}

	The key statements of Faltings are collected below for convenience.
\end{ups}

\begin{thm}[{\cite[Lemma~1]{Faltings1983}}]
	Let $S$ be a normal scheme, $U\sbs S$ open and dense, $p_{i}:A_{i}\ar S, i=1,2$ are two semiabelian varieties. If $\phi:A_{1}\tms_{S}U\ar A_{2}\tms_{S}U$ is a morphism of algebraic groups, then $\phi$ can be extended to $\oL{\phi}:A_{1}\ar A_{2}$.
\end{thm}

% \begin{thm}[{\cite[Lemma~2]{Faltings1983}}]
% 	Let $\kA_{g}$ be the algebraic stack of principally polarized abelian varieties (which is not over $\Spec(\bZ)$) and $p:\cA\ar\kA_{g}$ the universal abelian variety, $\cA_{g}\sbs\bP_{\bQ}^{N}$ given by $\oga_{\cA/\kA_{g}}^{\ot r}$ for some $r>0$. Then there exists a proper algebraic stack $\kZ$ over $\Spec(\bZ)$ and open substack $\kU\sbs\kZ$, and a proper morphism $\psi:\kU\ar\kA_{g}$, which extends to $\oL{\psi}:\kZ\ar\oL{\cA}_{g}$ over $\bQ$, so that the following objects exists
% 	\begin{enr}
% 		\item A universal stable curve $q:\kC\ar\kZ$.
% 		\item A line bundle $\cL\sbs\Lda^{g}q\zS(\oga_{\kC/\bZ})$ such that $\cL^{\ot r}\cong\oL{\psi}\xS(\cO(1))$ over $\kZ\ot\bQ$ and $\cL|_{\kU}^{\ot r}\cong\psi\xS(\cO(1))$ over $\kU\ot\bQ$.
% 		\item A pair of homomorphisms over $\kU$,
% 		\[
% 		  \afa:\Pic^{\tau}(\kC/\kZ)\ar\psi\xS(\cA),\quad \bta:\psi\xS(\cA)\ar\Pic^{\tau}(\kC/\kZ),
% 		\]
% 		such that $\afa\cc\bta=[d], d\in\bZ_{>0}$.
% 	\end{enr}
% \end{thm}

% \begin{thm}[{\cite[Lemma~3]{Faltings1983}}]
% 	Let $Y\sbs X\sbs\bP_{\bZ}^{n}$ be Zariski closed sets, $\dP{\cdot}$ the Hermitian metric on $\cO(1)|_{X(\bC)-Y(\bC)}$ with logarithmic singularities along $Y$. The for any number field $K$ and $c>0$, there are only finitely many $x\in X(K)-Y(K)$ with $h_{\cO(1)}(x)\lqs c$.
% \end{thm}

\begin{thm}[{\cite[Lemma~4]{Faltings1983}, Hermite-Minkowski}]
	Let $K$ be a number field, $S$ a finite set of finite places of $K$. Then there are only finitely many field extensions $K'\sps K$ of given order which is unramified outside $S$.
\end{thm}

\begin{thm}[{\cite[Lemma~5]{Faltings1983}, Faltings isogeny lemma}]
	Let $\phi:A_{1}\ar A_{2}$ be an isogeny of abelian varieties over a number field $K$ extending to an isogeny of N{\'e}ron models $\phi:\cA_{1}\ar\cA_{2}$ over $\cO_{K}$ with $\cG=\Ker\phi$. Then
	\[
		h_{\text{Fal}}(A_{2})-h_{\text{Fal}}(A_{1})=\frac{1}{2}\log(\Deg(\phi))-\frac{1}{[K:\bQ]}\log(\# s\xS\Oga_{\cG/\cO_{K}}^{1}).
	\]
\end{thm}

\begin{thm}[{\cite[Satz~1]{Faltings1983}}]
	Given $c>0$ and an abelian variety $A$ over $K$, then there are finitely many isomorphism classes of pairs $(\cA,\lda)$, where
	\begin{itm}
		\item $\cA\ar\Spec(\cO_{K})$ is an semiabelian variety of relative dimension $g$ and proper generic fiber $A$.
		\item $\lda$ is a principal polarization with $h_{\lda}(A)\lqs c$.
	\end{itm}
\end{thm}

\begin{thm}[{\cite[Satz~2]{Faltings1983,Faltings1983Erratum}}]
	Let $p:\cA\ar\Spec(\cO_{K})$ be a semiabelian variety with proper generic fiber, $\ell$ a rational prime and $G\sbs \cA[\ell\xF]$ an $\ell$-divisible subgroup over $K$. Then for $n\gg0$, $h_{\text{Fal}}(\cA/G_{n})=h_{\text{Fal}}(\cA)$.
\end{thm}

\begin{thm}[{\cite[Satz~3,4]{Faltings1983}, Tate conjectures}]
	Let $K$ be a number field, $A$ an abelian variety of dimension $g$ over $K$, $\ell$ a rational prime, $T_{\ell}(A)$ the Tate module with $G_{K}$-action. Then
	\begin{itm}
		\item The rational Tate module $V_{\ell}(A)=T_{\ell}(A)\ot_{\bZ_{\ell}}\bQ_{\ell}$ is semisimple.
		\item The natural map $\End_{K}(A)\ot_{\bZ}\bZ_{\ell}\ar\End_{G_{K}}(T_{\ell}(A))$ is an isomorphism.
	\end{itm}

	Moreover, let $A_{1},A_{2}$ be abelian varieties over $K$, then
	\[
		\Hom_{K}(A_{1},A_{2})\ot_{\bZ}\bZ_{\ell}\cong\Hom_{G_{K}}(T_{\ell}(A_{1}),T_{\ell}(A_{2})),
	\]
	and the following are equivalent
	\begin{enr}
		\item $A_{1}$ and $A_{2}$ isogenous.
		\item $V_{\ell}(A_{1})\cong V_{\ell}(A_{2})$ as $G_{K}$-modules.
		\item $L_{v}(s,A_{1})=L_{v}(s,A_{2})$ for almost all finite places $v$ of $K$.
		\item $L_{v}(s,A_{1})=L_{v}(s,A_{2})$ for all finite places $v$ of $K$.
	\end{enr}

	% Let $A$ be an abelian variety, then there are finitely many isomorphism classes of abelian varieties $B$ of polarization degree $d$ with $T_{\ell}(A)\cong T_{\ell}(B)$ for all $p$.
\end{thm}

\begin{thm}[{\cite[Satz~5,6]{Faltings1983}, Shafarevich conjectures}]
	Let $S$ be a finite set of finite places of $K$, then there are finitely many
	\begin{enr}
		\item isogeny classes of abelian varieties of given dimension,
		\item isomorphism classes of abelian varieties of given polarization degree and dimension,
		\item isomorphism classes of smooth curves of genus $g\gqs2$,
	\end{enr}
	over $K$ which have good reduction outside $S$.
\end{thm}

\begin{thm}[{\cite[Satz~7]{Faltings1983}, Mordell conjecture}]
	Let $X$ be a smooth curve of genus $g\gqs2$ over $K$, then $X(K)$ is finite.
\end{thm}




\ssc{CO2009}

\begin{stp}
	Let $\cA_{g,d,n}\ar\Spec(\bZ[1/dn])$ be the coarse moduli scheme of polarized abelian varieties of dimension $g$ and polarization degree $d$ with a symplectic level-$n$ structure. Let $\cA_{g,d}\ar\Spec(\bZ)$ be the coarse moduli scheme of polarized abelian varieties of dimension $g$ and polarization degree $d$. Let $\cA_{g}=\coprod_{d}\cA_{g,d}\ar\Spec(\bZ)$ be the coarse moduli scheme of polarized abelian variety of dimension $g$. Let $K\sps\bF_{p}$ be a perfect field throughout this subsection unless otherwise specified.
\end{stp}

\prg{Main results}

\begin{thm}[{\cite[1.8, 1.13]{CO2009}}]
	Let $[(A,\lda)]=x\in\cA_{g}\tms_{\Spec(\bZ)}\Spec(\bF_{p})$ be the moduli point of a polarized abelian variety, then
	\begin{enr}
			\item (Chai) If $A$ is ordinary, then $\cH(x)$ is Zariski dense in $\cA_{g}\tms_{\Spec(\bZ)}\Spec(\bF_{p})$.
			\item (Chai-Oort) $\cH(x)$ is Zariski dense in $\cW_{\cN(A)}(\cA_{g}\tms_{\Spec(\bZ)}\Spec(\bF_{p}))$.
	\end{enr}
\end{thm}

\begin{thm}[{\cite[1.19, 1.20, 1.22, 7.3]{CO2009}}]
	One has
	\begin{enr}
			\item (Grothendieck-Katz) Let $S$ be an integral scheme over $K$, $\cX\ar S$ a $p$-divisible group, then given any Newton polygon $\zta$, $\cW_{\zta}(S)=\bmat{s\in S: \cN(\cX_{s})\prc \zta}\sbs S$ is closed. Let $\eta$ be the generic point of $S$, if $\zta=\cN(\cX_{\eta})$, then $D=\bmat{s\in S:\cN(\cX_{s})\pcn \zta}\sbs S$ is closed, and is either empty or of pure codimension $1$.
			\item (Grothendieck-Oort) Given any $p$-divisible group $X\ar\Spec(K)$ and any Newton polygon $\cN(X)\prc \gma$, there exists an integral scheme $S$ over $K$ with generic point $\eta$, a specialization $s$ of $\eta$, and a $p$-divisible group $\cX\ar S$ such that $\cN(X_{\eta})=\gma$ and $\cX_{s}\cong X$.
	\end{enr}
\end{thm}

\begin{rmk}
	The ordinary Hecke orbits are dense, and every $p$-divisible group can be realized as the deformation of another $p$-divisible group with lower Newton polygon.
\end{rmk}


\prg{Serre-Tate theory}

\begin{ups}
	Deforming an abelian variety over $K$ and deforming its $p$-divisible group is the same (in the sense of categorical equivalence), and the deformation functor of the latter is representable by a smooth formal scheme. The deformation functor of ordinary $p$-divisible groups has a structure of formal torus.
\end{ups}

\begin{dfn}
	Let $S_{0}$ be a scheme over $K$, the category of deformations $\Def(S_{0}/W(K))$ of $S_{0}$ contains objects of pairs $(S\ar\Spec(R),\eps)$, where $R$ is an Artinian local algebra over $W(K)$ and $\eps$ is an isomorphism,
	\[
		\begin{tikzcd}
			S\ar[d] & S\tms_{\Spec(R)}\Spec(R/\km_{R})\ar[r, equal, "\eps"]\ar[l]\ar[d] & S_{0}\tms_{\Spec(K)}\Spec(R/\km_{R})\ar[r]\ar[d] & S_{0}\ar[d] \\
			\Spec(R) & \Spec(R/\km_{R})\ar[l]\ar[r, equal] & \Spec(R/\km_{R})\ar[r] & \Spec(K)
		\end{tikzcd}
	\]
\end{dfn}

\begin{thm}[{\cite[2.5, 2.8]{CO2009}}]
	Let $X_{0}$ be a $p$-divisible group and $A_{0}$ be an abelian variety over $K$, then
	\begin{enr}
		\item (Grothendieck-Messing) $\Def(X_{0}/W(K)):\Art_{W(K)}\ar\Set$ is representable,
		\[
			\Def(X_{0}/W(K))\cong\Spf(W(K)\dS{x_{1},\ldots,x_{\Dim(X_{0})\Dim(X_{0}\xV)}})
		\]
		is a non-canonical isomorphism.
		\item (Serre-Tate) $(A\ar\Spec(R),\eps)\amt(A[p\xF],\eps[p\xF])$ gives an equivalence of categories,
		\[
			\bG:\Def(A_{0}/W(K))\cong\Def(A_{0}[p\xF]/W(K)),
		\]
		moreover, $\Def(A_{0}/W(K))$ admits an action by $\Aut(A_{0}[p\xF])$.
	\end{enr}
\end{thm}

\begin{thm}[{\cite[2.19]{CO2009}}]
	Let $X_{0}$ be an ordinary $p$-divisible group over $K$ with $0\ar T_{0}\ar X_{0}\ar E_{0}\ar 0$, lift the toric and {\'e}tale part to $\oT{T}$ and $\oT{E}$ over $\Spec(W(K))$, and $\oH{T}$ the formal torus of $T_{0}$ over $\Spec(W(K))$, then
	\begin{enr}
		\item Every $(X\ar\Spec(R),\eps)\in\Def(X_{0}/W(K))$ is ordinary with
		\[
			0\ar\oT{T}\tms_{\Spec(W(K))}\Spec(R)\ar X\ar \oT{E}\tms_{\Spec(W(K))}\Spec(R)\ar 0,
		\]
		therefore $\Def(X_{0}/W(K))$ can be defined as a functor $\Art_{W(K)}\ar\Ab$ via Baer sum construction.
		\item There is a natural isomorphism of functors,
		\[
			\uL{\Hom}_{\bZ_{p}}(T_{p}(E_{0}),\oH{T})\cong\Def(X_{0}/W(K)),
		\]
		therefore $\Def(X_{0}/W(K))$ is a formal torus whose cocharacter group is $T_{p}(E)\xV\ot_{\bZ_{p}}X\zS(T_{0})$.
	\end{enr}
\end{thm}

\prg{Tate conjecture}
\begin{ups}
	The Tate conjecture of abelian varieties is true over finite fields by Tate and number fields by Faltings.
\end{ups}

\begin{thm}[{\cite[3.3, 3.9, 3.12]{CO2009}}]
	Let $A$ be a simple abelian variety over $\bF_{q}$, then
		\begin{enr}
			\item (Tate) There is a CM-field of degree $2\Dim(A)$ contained in $\End(A)_{\bQ}$.
			\item (Weil) $F^{(n)}:A\ar A^{(p^{n})}=A$ determines a number $\pi_{A}\in\End(A)_{\bQ}$, $\pi_{A}$ is a Weil $q$-number.
		\end{enr}
	(Honda-Tate) Moreover, the isogeny classes of simple abelian varieties over $\bF_{q}$ is bijective to the conjugacy classes of Weil $q$-numbers.
\end{thm}

\begin{thm}[{\cite[3.16, 3.20]{CO2009}}]
	Let $A$ be an abelian variety over a finitely generated field $K$, then
	\begin{enr}
			\item (Tate-Faltings) If $\ell\neq\Char(K)$, then there is an isomorphism of $G_{K}$-modules,
		\[
			\End(A)_{\bZ_{\ell}}\cong\End(T_{\ell}(A)).
		\]
			\item (Tate-de Jong) If $p=\Char(K)>0$, then there is an isomorphism of $G_{K}$-modules,
		\[
			\End(A)_{\bZ_{p}}\cong\End(A[p\xF],B[p\xF]).
		\]
	\end{enr}
\end{thm}


\prg{Cartier-Dieudonn{\'e} theory}

\begin{ups}
	The theory of Cartier modules (resp., Dieudonn{\'e} modules) give a linear algebric description (in the sense of categorical equivalence) of commutative smooth formal groups (resp., $p$-divisible formal groups) in terms of modules over suitable non-commutative rings.
\end{ups}

\begin{stp}
	Denote by $W(K)$ the ring of Witt vectors of $K$, $L=\Frac(W(K))$ and $\sgm:W(K)\ar W(K)$ the Teichm{\"u}ller lift of the Frobenius on $K$, denote by $F,V$ the Frobenius and Verschiebung operators.
\end{stp}

\begin{dfn}[{\cite[4.28]{CO2009}}]
	Define $R_{K}$ and $\Cart_{p}(K)$ as
	\begin{align*}
	  R_{K} &= \bmat{\sum_{i\in\bZ}a_{i}V^{i}: \fal i\in\bZ, a_{i}\in L, \ord_{p}(a_{i})\gqs\max(0,-i), \text{ and $a_{i}=0$ for all but finitely many $i$}},\\
	  \Cart_{p}(K) &= \bmat{\sum_{i\in\bZ}a_{i}V^{i}: \fal i\in\bZ, a_{i}\in L, \ord_{p}(a_{i})\gqs\max(0,-i), \text{ and $\ord_{p}(a_{i})+i\ar\ift$ as $i\ar\ift$}},
	\end{align*}
	where $x\cdot V^{i}=V^{i}\cdot\sgm^{i}(x), \fal x\in W(K), i\in\bZ$. A Cartier module is a left $\Cart_{p}(K)$-module $M$. A Dieudonn{\'e} module is a left $R_{K}$-module $M$ which is free as $W(K)$-module of finite rank.
\end{dfn}


\begin{thm}[{\cite[4.31, 4.33, 4.37, 4.42]{CO2009}}]
	Let $X$ be a $p$-divisible group over $K$, then
	\begin{enr}
			\item There is a canonical splitting over $K$, $X\cong X_{\text{to}}\tms X_{\ell\ell}\tms X_{\text{{\'e}t}}$, where $X_{\text{{\'e}t}}$ is the maximal {\'e}tale quotient, $X_{\text{tor}}$ is the maximal toric sub, and $X_{\ell\ell}$ admits neither non-trivial {\'e}tale quotient nor toric sub, then
		\begin{enr}
				\item $X=X_{\text{{\'e}t}}$ if and only if $V$ is bijective or $F$ is divisible by $p$.
				\item $X=X_{\text{to}}$ if and only if $F$ is bijective or $V$ is divisible by $p$.
				\item $X=X_{\ell\ell}$ if and only if $F,V$ are both topologically nilpotent.
		\end{enr}
			\item (Dieudonn{\'e}-Manin Classification) There is an isogeny over $\oL{K}$, $X\atr G_{m_{1},n_{1}}\tms\cdots\tms G_{m_{r},n_{r}}$, satisfying for all $1\lqs i\lqs r$, $\bD(G_{m_{i},n_{i}})=R_{K}/(V^{n_{i}}-F^{m_{i}}), \Gcd(m_{i},n_{i})=1$, and for all $1\lqs i\lqs r-1$, $m_{i}/(m_{i}+n_{i})\lqs m_{i+1}/(m_{i+1}+n_{i+1})$ determines its Newton polygon $\cN(X)$, in fact, the isogeny class of $X$ over $\oL{K}$ is classified by $\cN(X)$, in particular,
		\begin{enr}
				\item $X$ is {\'e}tale if and only if all its slopes are $0$.
				\item $X$ is toric if and only if all its slopes are $1$.
		\end{enr}
			\item The category of $p$-divisible groups over $K$ and the category of Dieudonn{\'e} modules over $R_{K}$ are equivalent via $X\amt\bD(X)$, compatible with direct product and exactness.
			\item If $X_{\text{{\'e}t}}$ is trivial, let $\oH{X}$ be the formal completion of $X$ along the zero section of $X$. Then there is a canonical isomorphism between Dieudonn{\'e} module and Cartier module $\bD(X)\cong\bM_{p}(\oH{X})$, compatible with the action by $F,V,W(K)$.
			\item One has $\Ht(X)=\Rnk_{W(K)}(\bD(X))$ and $\Lie(X)\cong\bD(X)/V\bD(X)$.
			\item One has $\bD(X\xV)\cong\bD(X)\xV=\Hom_{W(K)}(\bD(X),W(K))$ as $W(K)$-modules, and action of $F,V$ given by
		\[
			(V\cdot\vfi)(m)=\sgm\xI(h(Fm)),\quad(F\cdot\vfi)(m)=\sgm(h(Vm)),\quad\fal h\in\bD(X)\xV, m\in\bD(X).
		\]
	\end{enr}
\end{thm}

% \ssc{MadapusiPera2019}



\ssc{TODO~arXiv:2310.14428}
This paper gives some relations between height of modular polynomials and height of $X_{0}(N)$, that is $h_{\mathrm{Fal}}(X_{0}(N))\sim h(\Phi_{N})/6^{3}$.




% ----------------------------------------------------------------------------------------------------
\bibliographystyle{amsalpha}
\bibliography{/Users/wenc/Documents/Math/universal-ref}
\end{document}
